\documentclass[12pt,a4paper]{article}
% Drake_preamble.tex - LaTeX Preamble Template
% 作者: 謝慕揚, MD, PhD, FESC
% 
% 使用說明:
% 1. 表格請使用 longtable，不要有垂直分隔線 |
% 2. 表格內換行使用 \newline，不要使用 \\
% 3. 藥物名稱、檢驗、檢查請維持英文
% 4. Reference 格式請使用 \href 加上驗證過的 DOI

% Including essential packages for Chinese typesetting
\usepackage{xeCJK}
\usepackage{fontspec}
% Configuring Chinese font with Noto Serif CJK TC, fallback to SimSun
\IfFontExistsTF{Noto Serif CJK TC}{
  \setCJKmainfont{Noto Serif CJK TC}
  \setCJKsansfont{Noto Serif CJK TC}
  \setCJKmonofont{Noto Serif CJK TC}
}{\setCJKmainfont{SimSun}
  \setCJKsansfont{SimSun}
  \setCJKmonofont{SimSun}
}

% Including other necessary packages
\usepackage{geometry}
\usepackage{titlesec}
\usepackage{enumitem}
\usepackage{xcolor}
\usepackage{colortbl}
\usepackage{tcolorbox}
\usepackage{booktabs}
\usepackage{longtable}
\usepackage{array}
\usepackage{graphicx}
\usepackage{tikz}
\usepackage{fancyhdr}
\usepackage{multicol}
\usepackage{datetime2}
\usepackage{hyperref}
\usepackage{mdframed}
\usepackage{amsmath}
\usepackage{amssymb}
\usepackage{multirow}

\hypersetup{
    colorlinks=true,
    linkcolor=emergency_blue,
    citecolor=emergency_blue,
    urlcolor=emergency_blue,
    pdfborder={0 0 0}
}

% Defining page geometry
\geometry{
  top=2.5cm,
  bottom=2.5cm,
  left=2cm,
  right=2cm,
  headheight=25pt
}

% Adjusting topmargin to compensate for increased headheight
\addtolength{\topmargin}{-10pt}

% Defining colors
\definecolor{emergency_red}{RGB}{186,24,27}
\definecolor{emergency_blue}{RGB}{0,114,188}
\definecolor{emergency_gray}{RGB}{120,120,120}
\definecolor{highlight_yellow}{RGB}{255,250,205}

% Setting title formats
\titleformat{\section}[block]{\Large\bfseries\color{emergency_red}}{\thesection}{10pt}{}
\titleformat{\subsection}[block]{\large\bfseries\color{emergency_blue}}{\thesubsection}{8pt}{}
\titleformat{\subsubsection}[block]{\normalsize\bfseries\color{emergency_gray}}{\thesubsubsection}{6pt}{}

% Defining custom environment for important notes
\newmdenv[
    linecolor=emergency_red,
    backgroundcolor=highlight_yellow,
    linewidth=2pt,
    topline=true,
    bottomline=true,
    leftline=true,
    rightline=true,
    innertopmargin=10pt,
    innerbottommargin=10pt,
    innerleftmargin=10pt,
    innerrightmargin=10pt
]{importantbox}

% Defining note command with tcolorbox
\newcommand{\note}[1]{
    \par
    \noindent
    \begin{tcolorbox}[
        colback=emergency_blue!5!white,
        colframe=emergency_blue,
        title=\textbf{重點提醒},
        fonttitle=\bfseries,
        boxrule=0.5pt,
        arc=3pt,
        left=6pt,
        right=6pt,
        top=6pt,
        bottom=6pt
    ]
    #1
    \end{tcolorbox}
}

% Key findings box
\newcommand{\keyfinding}[1]{
    \par
    \noindent
    \begin{tcolorbox}[
        colback=yellow!10!white,
        colframe=orange!75!black,
        title=\textbf{核心發現摘要},
        fonttitle=\bfseries,
        boxrule=0.5pt,
        arc=3pt
    ]
    #1
    \end{tcolorbox}
}

% Defining custom date format for ISO 8601
\DTMnewdatestyle{iso}{
  \renewcommand{\DTMdisplaydate}[4]{\number##1-\number##2-\number##3}
  \renewcommand{\DTMDisplaydate}[4]{\number##1-\number##2-\number##3}
}
\DTMsetdatestyle{iso}
\newcommand{\mydate}{\DTMtoday}


% Custom header for this document
\pagestyle{fancy}
\fancyhf{}
\fancyhead[L]{\textcolor{emergency_red}{JACC 教學講義}}
\fancyhead[R]{\textcolor{emergency_blue}{COVID-19 與心血管預後}}
\fancyfoot[C]{\textcolor{emergency_gray}{\thepage}}
\renewcommand{\headrulewidth}{1pt}
\renewcommand{\headrule}{\hbox to\headwidth{\color{emergency_red}\leaders\hrule height \headrulewidth\hfill}}

\begin{document}

%============================================================
% Title Page
%============================================================
\begin{center}
{\LARGE\bfseries\textcolor{emergency_red}{Journal of the American College of Cardiology}}\\[0.5cm]
{\Large\textbf{Hospitalizations and Mortality Among Medicare Beneficiaries at Cardiovascular Risk: Disruptions and Recovery from the COVID-19 Pandemic}}\\[0.3cm]
{\large 心血管風險 Medicare 受益人之住院率與死亡率：COVID-19 疫情的衝擊與恢復}\\[0.5cm]
{\normalsize \href{https://www.jacc.org}{JACC 2026}}\\[0.3cm]
{\normalsize 整理：謝慕揚 MD, PhD, FESC \quad 日期：\mydate}
\end{center}

\vspace{0.5cm}

%============================================================
\section{研究背景}
%============================================================

COVID-19 疫情對心血管照護造成重大衝擊。2020 年疫情高峰期間，急性心血管疾病住院率大幅下降的現象已有文獻記載。然而，疫情「後期」（2021-2022 年）的醫療利用與病患預後變化，尚缺乏系統性的研究。

\subsection{研究目的}
分析疫情後期（2021-2022 年）Medicare 受益人中，具心血管風險因子或心血管疾病者的：
\begin{itemize}
    \item 醫療利用模式（住院、門診）
    \item 病患預後（死亡率）
\end{itemize}

%============================================================
\section{研究方法}
%============================================================

\begin{itemize}
    \item \textbf{研究設計}：流行病學回顧性世代研究
    \item \textbf{資料來源}：Medicare 資料庫（Fee-for-Service 與 Medicare Advantage）
    \item \textbf{研究族群}：約 5,000 萬名 Medicare 受益人
    \item \textbf{比較時期}：
    \begin{itemize}
        \item 疫情前：2018-2019 年
        \item 疫情後期：2021-2022 年
    \end{itemize}
    \item \textbf{納入條件}：具心血管風險因子或心血管疾病
\end{itemize}

%============================================================
\section{主要發現}
%============================================================

\keyfinding{
\begin{itemize}
    \item 住院率較疫情前下降 \textbf{8-9\%}（未完全恢復）
    \item 門診就醫率恢復並超越基線（受 telehealth 推升）
    \item 死亡率持續升高約 \textbf{20\%}（相對增加）
\end{itemize}
}

\subsection{醫療利用趨勢}

\begin{longtable}{p{3.5cm}p{10.5cm}}
\toprule
\textbf{指標} & \textbf{變化趨勢} \\
\midrule
\endfirsthead
\toprule
\textbf{指標} & \textbf{變化趨勢} \\
\midrule
\endhead
\bottomrule
\endfoot

住院率 & 疫情前穩定 → 2020 年驟降 → 恢復但未達疫情前水準\newline （較疫情前低 8-9\%） \\
\midrule
門診就醫率 & 疫情前穩定 → 2020 年驟降 → 恢復並超越基線\newline （Telehealth 的興起為主因） \\
\midrule
死亡率 & 疫情前穩定趨勢 → 2020 年急升 → 持續波動升高\newline （隨 COVID-19 波次起伏，整體較疫情前高約 20\%） \\
\bottomrule
\end{longtable}

%============================================================
\section{次族群分析}
%============================================================

\subsection{住院率下降最明顯的族群}
\begin{itemize}
    \item 鄉村地區（rural areas）
    \item 社會脆弱度高的地區（high Social Vulnerability Index, SVI）
    \item Medicare Advantage 受益人（下降幅度較 Fee-for-Service 小）
\end{itemize}

\subsection{門診就醫率上升最明顯的族群}
\begin{itemize}
    \item 都市地區（urban areas）
    \item Medicare Advantage 受益人
\end{itemize}

\subsection{死亡率上升最明顯的族群}

\begin{importantbox}
\textbf{高風險族群}
\begin{itemize}
    \item 社會脆弱度高的族群（high Social Vulnerability Index）
    \item Medicare Advantage 受益人（較 Fee-for-Service 更明顯）
\end{itemize}

\textbf{注意}：鄉村 vs 都市地區的死亡率增幅\textbf{無顯著差異}
\end{importantbox}

%============================================================
\section{可能機轉探討}
%============================================================

作者提出三大可能機轉：

\subsection{1. 生物學因素（Biological）}
\begin{itemize}
    \item COVID-19 感染本身的急性危害
    \item COVID-19 感染後的長期影響（long COVID）
    \item 瑞典研究顯示：COVID-19 感染者在一年後仍有較高的全因死亡率與心血管死亡風險
    \item 呼吸道病毒（流感、COVID-19、副流感病毒等）可觸發心衰竭、缺血性中風、急性心肌梗塞
\end{itemize}

\subsection{2. 行為因素（Behavioral）}
\begin{itemize}
    \item 疫情期間的就醫迴避行為（care avoidance）持續存在
    \item 影響較大的族群：
    \begin{itemize}
        \item 女性
        \item 部分種族/族裔少數群體（原本就醫可近性較差）
    \end{itemize}
    \item 處置積壓（backlogs for procedures）
    \item 預防保健延遲
\end{itemize}

\subsection{3. 醫療系統因素（Systems）}
\begin{itemize}
    \item 醫院人力吃緊（hospital strain）
    \item 照護模式改變
    \item 技術性護理機構（Skilled Nursing Facility, SNF）供給減少
    \begin{itemize}
        \item 人力短缺問題
        \item 患者住院天數延長
        \item 可能影響功能狀態與死亡率
    \end{itemize}
\end{itemize}

%============================================================
\section{臨床與公衛意涵}
%============================================================

\note{
\textbf{為下一次公衛緊急事件做準備}

本研究揭示了心血管照護系統的脆弱點：
\begin{itemize}
    \item 社會脆弱族群受影響最大
    \item 醫療利用與死亡率的解離（住院減少但死亡增加）
    \item 需強化醫療系統韌性（resilience）
\end{itemize}
}

\subsection{值得關注的議題}
\begin{enumerate}
    \item 醫療利用下降是否反映「適當減少不必要住院」或「必要照護的延誤」？
    \item 死亡率持續升高的驅動因素為何？
    \item 如何縮小社會脆弱族群的健康差距？
    \item Medicare Advantage 與 Fee-for-Service 的預後差異原因？
\end{enumerate}

%============================================================
\section{重點摘要}
%============================================================

\begin{enumerate}
    \item \textbf{住院率}：疫情後期仍較疫情前低 8-9\%，未完全恢復
    \item \textbf{門診就醫}：恢復並超越基線，Telehealth 為主要推手
    \item \textbf{死亡率}：持續較疫情前高約 20\%，令人憂心
    \item \textbf{高風險族群}：社會脆弱度高者、Medicare Advantage 受益人
    \item \textbf{多重機轉}：生物學、行為、系統因素交互作用
    \item \textbf{公衛啟示}：需強化醫療系統韌性，為未來緊急事件做準備
\end{enumerate}

%============================================================
\section{參考文獻}
%============================================================

\begin{enumerate}
    \item Maddox KJ, Wadhera R, Waken A, et al. Hospitalizations and Mortality Among Medicare Beneficiaries at Cardiovascular Risk: Disruptions and Recovery from the COVID-19 Pandemic. \textit{J Am Coll Cardiol}. 2026.

    \item Wadhera R, Shen C, Gondi S, et al. Cardiovascular Deaths During the COVID-19 Pandemic in the United States. \href{https://doi.org/10.1016/j.jacc.2020.10.055}{\textit{J Am Coll Cardiol}. 2021;77(2):159-169.}

    \item Solomon MD, McNulty EJ, Rana JS, et al. The Covid-19 Pandemic and the Incidence of Acute Myocardial Infarction. \href{https://doi.org/10.1056/NEJMc2015630}{\textit{N Engl J Med}. 2020;383:691-693.}

    \item Wadhera R, Dhruva S, Bikdeli B, et al. Cardiovascular Statistics in the United States, 2026: JACC Stats. \href{https://doi.org/10.1016/j.jacc.2025.12.027}{\textit{J Am Coll Cardiol}. 2026.}
\end{enumerate}

\end{document}
